%%%%%%%%%%%%%%%%%%%%%%%%%%%%%%%%%%%%%%%%%%%%%%%%%%%%%%%%%%%%%%%%%%%%%%%%%%%
%%  chapter3.tex  –  Глава 3: Адаптация и согласование с системами
%%                   обнаружения и предотвращения сетевых вторжений
%%  v9 – Предметно-ориентированная глава: терминологии, преобразование
%%       данные→граф, таксономия (34 класса), признаки (83), гиперпараметры,
%%       реализации методов. Без экспериментальных результатов (см. Главу 4).
%%       Только заголовки \section добавлены в оглавление.
%%%%%%%%%%%%%%%%%%%%%%%%%%%%%%%%%%%%%%%%%%%%%%%%%%%%%%%%%%%%%%%%%%%%%%%%%%%

\chapter{Адаптация и согласование с системами обнаружения и предотвращения сетевых вторжений}
\label{ch:adaptation}

В настоящей главе предметно-независимые математические методы, разработанные в Главе~\ref{ch:methods}, отображаются на сущности и ограничения кибербезопасности и, в частности, гибридных систем обнаружения и предотвращения сетевых вторжений (СОПСВ). Абстрактные объекты предыдущей главы --- темпоральные графы, динамика пространства состояний, вариационные апостериорные распределения, транспортные планы и теоретико-игровые полезности --- реализуются с конкретными сущностями сетевой безопасности: потоками, пакетами, сессиями, хостами, сервисами и злонамеренными субъектами. В главе определены специфические для кибербезопасности терминологии, описано преобразование сырого сетевого трафика в графовые структуры, пригодные для ввода в разработанные архитектуры ГНС, определены таксономия классификации и конвейер инженерии признаков, представлена конфигурация гиперпараметров и обоснование оптимизации для каждой модели, продемонстрировано, как каждый метод из Главы~\ref{ch:methods} адаптирован с предметно-ориентированными модификациями. Экспериментальные результаты здесь не представлены; они зарезервированы для Главы~\ref{ch:experiments}.


%%=====================================================================
\section{Обнаружение сетевых вторжений: терминологии и контекст}
\label{sec:nidps_context}
%%=====================================================================

\subsection{Гибридные системы обнаружения и предотвращения вторжений}

Гибридная система обнаружения и предотвращения вторжений (ГСОПВ)~\cite{Buczak2016,Ahmad2021,Aldweesh2020,Ferrag2020,Paxson1999} объединяет две взаимодополняющие парадигмы обнаружения. Сигнатурное обнаружение сопоставляет входящий трафик с базой данных известных паттернов (правил) и достигает высокой прецизионности на ранее каталогизированных угрозах, но не может обнаруживать новые атаки. Аномальное обнаружение обучает статистическую модель нормального поведения и маркирует отклонения, обеспечивая обнаружение эксплойтов нулевого дня и эволюционирующих паттернов атак ценой потенциально более высокого уровня ложных срабатываний. Методы диссертации работают в рамках аномальной ветви, дополненной рассуждением большой языковой модели для классификации zero-shot новых категорий.

Сетевая СОПВ (ССОПВ) мониторит трафик на сетевом уровне, анализируя потоки и пакеты между хостами, а не события внутри отдельного конечного устройства. Все три рабочих условия архитектурной основы (Рис.~\ref{fig:dissertation_architecture}) присутствуют в данной области: состязательная манипуляция целенаправленными противниками, формирующими трафик для уклонения от обнаружения, нестационарность через эволюционирующие паттерны трафика и изменения протоколов, и распределённость инфраструктуры мониторинга по организационным границам.


\subsection{Ключевые терминологии}

Следующие термины используются на протяжении данной главы и экспериментальной валидации в Главе~\ref{ch:experiments}:

\noindent\textbf{Поток.} Последовательность пакетов с общей 5-кортежной характеристикой (IP-адрес источника, IP-адрес назначения, порт источника, порт назначения, протокол). Поток является фундаментальной единицей анализа и соответствует направленному ребру в сетевом графе.

\noindent\textbf{Сессия.} Двунаправленная коммуникация, включающая прямой поток и его обратный аналог, захватывающая полный обмен между двумя конечными точками.

\noindent\textbf{Хост.} Сетевая конечная точка (IP-адрес), участвующая как источник или назначение. Хосты соответствуют узлам в сетевом графе.

\noindent\textbf{Сервис.} Функциональный компонент, идентифицируемый номером порта и протоколом (например, HTTP на порту 80, SSH на порту 22). В облачных развёртываниях сервисы соответствуют контейнерам микросервисов.

\noindent\textbf{PCAP (Packet Capture).} Сырой двоичный формат, захватывающий полные заголовки пакетов и опционально полезные нагрузки, из которого извлекаются признаки уровня потоков.

\noindent\textbf{Правило СОВ.} Структурированная спецификация паттерна (например, в синтаксисе правил Snort или Suricata), кодирующая сигнатуру обнаружения с действием, протоколом, адресами, портами и директивами сопоставления содержимого.

\noindent\textbf{MITRE ATT\&CK.} Комплексная база знаний о тактиках, техниках и процедурах (ТТП) противников, используемая для классификации и контекстуализации обнаруженных угроз.

\noindent\textbf{SOC (Центр операций безопасности).} Организационное подразделение, ответственное за мониторинг, анализ и реагирование на события безопасности, укомплектованное аналитиками уровней~1--3.


%%=====================================================================
\section{Конвейер преобразования данных в граф}
\label{sec:data_to_graph}
%%=====================================================================

Центральная предпосылка настоящей диссертации состоит в том, что сетевой трафик имеет естественную графовую структуру: хосты являются узлами, потоки --- рёбрами, а временной порядок определяет эволюцию графа. Данный раздел формализует преобразование сырых сетевых данных в графовые входы, требуемые разработанными архитектурами ГНС.

\subsection{Построение графа из сетевого трафика}

При заданном потоке сетевых записей $\{(s_i, d_i, t_i, \mathbf{f}_i)\}_{i=1}^{N}$, где $s_i$ --- хост-источник, $d_i$ --- хост-назначение, $t_i$ --- временная метка и $\mathbf{f}_i\in\mathbb{R}^{83}$ --- вектор признаков, темпоральный граф $\calG(t)=(\calV,\calE(t),\mathbf{X}(t))$ строится следующим образом:

\noindent\textbf{Узлы.} Каждый уникальный IP-адрес становится узлом $v\in\calV$. Признаки узла $\bx_v(t)$ агрегируют статистику уровня хоста, вычисленную по скользящему окну $[t-\Delta_w, t]$: кумулятивные счётчики байтов (отправленных/полученных), счётчики активных соединений, уникальные порты назначения, распределение протоколов (соотношения TCP/UDP/ICMP) и средний межсоединительный интервал. Это даёт 12-мерный вектор признаков узла.

\noindent\textbf{Рёбра.} Каждый поток $(s_i, d_i, t_i, \mathbf{f}_i)$ создаёт направленное ребро от узла $s_i$ к узлу $d_i$ в момент $t_i$ с признаками ребра $\mathbf{e}_{s_id_i}(t_i) = \mathbf{f}_i$. Полный 83-мерный вектор признаков (Раздел~\ref{sec:features}) прикрепляется к ребру.

\noindent\textbf{Темпоральная эволюция.} Граф эволюционирует по мере поступления новых потоков и истечения старых соединений. Поток считается активным, если его последний пакет поступил в пределах настраиваемого тайм-аута $\Delta_{\mathrm{timeout}}=120$\,с. Истекшие рёбра удаляются, порождая динамическую, изменяющуюся во времени топологию.

\subsection{Многоуровневое построение графов}

Для TripleE-TGNN (Метод М3) и интеграции SOC-аналитики сетевой трафик дополнительно отображается на три иерархических уровня:

\noindent\textbf{Граф сервисов $\calG_s(t)$.} Узлы представляют сервисы (пары порт/протокол); рёбра представляют API-вызовы между сервисами, агрегированные по настраиваемым окнам. Используется для обнаружения паттернов межсервисных атак в средах микросервисов.

\noindent\textbf{Граф трассировок $\calG_r(t)$.} Узлы представляют распределённые трассировки запросов (идентифицируемые по ID трассировки в HTTP-заголовках или инструментации OpenTelemetry); рёбра соединяют трассировки, разделяющие хотя бы один сервис. Используется для обнаружения многоэтапных цепочек атак.

\noindent\textbf{Граф хостов $\calG_n(t)$.} Узлы представляют узлы инфраструктуры; рёбра представляют наблюдаемые сетевые потоки. Это основной граф для всех методов.

\noindent Двудольные графы связи $\calB_{sr}$, $\calB_{sn}$, $\calB_{rn}$ соединяют сущности между уровнями, обеспечивая механизм межуровневого внимания из уравнения~\eqref{eq:triple_cross_attn}.


%%=====================================================================
\section{Таксономия классификации}
\label{sec:taxonomy}
%%=====================================================================

Платформа классифицирует сетевой трафик в 34 различных класса: 1 легитимный класс и 33 класса атак, организованные в 7 основных категорий. Эта таксономия унифицирована по всем шести бенчмарк-наборам данных, при этом специфические для набора типы атак сопоставлены каноническим категориям. Таблица~\ref{tab:taxonomy} представляет полную таксономию.

\begin{table}[htbp]
\centering
\caption{Унифицированная таксономия классификации: 34 класса по 7 категориям атак.}
\label{tab:taxonomy}
\small
\begin{tabular}{p{2.5cm} p{10.5cm}}
\toprule
\textbf{Категория} & \textbf{Классы атак} \\
\midrule
\textbf{DDoS} (8) & DDoS-TCP, DDoS-UDP, DDoS-ICMP, DDoS-HTTP, DDoS-SYN Flood, DDoS-DNS Amplification, DDoS-SSDP, DDoS-NTP Amplification \\
\addlinespace
\textbf{DoS} (5) & DoS GoldenEye, DoS Hulk, DoS Slowhttptest, DoS Slowloris, DoS HeartBeat \\
\addlinespace
\textbf{Разведка} (4) & PortScan, Nmap Scan, Service Probe, OS Fingerprint \\
\addlinespace
\textbf{Перебор} (4) & FTP-Patator, SSH-Patator, Web Login Brute Force, RDP Brute Force \\
\addlinespace
\textbf{Веб-атаки} (4) & XSS, SQL Injection, Command Injection, Directory Traversal \\
\addlinespace
\textbf{ВПО/Ботнет} (5) & Mirai, Botnet Communication, Trojan Callback, Ransomware Beacon, Cryptomining \\
\addlinespace
\textbf{Продвинутые} (3) & Infiltration, Heartbleed, Spoofing \\
\addlinespace
\textbf{Легитимный} (1) & Нормальный трафик \\
\bottomrule
\end{tabular}
\end{table}

Таксономия расширяема: новые классы атак могут быть зарегистрированы через интерфейс администрирования платформы без переобучения полного ансамбля моделей, используя возможность классификации zero-shot CyberSecLLM и упругую консолидацию весов TripleE-TGNN для инкрементального обучения.


%%=====================================================================
\section{Конвейер инженерии признаков}
\label{sec:features}
%%=====================================================================

Все модели работают на общем векторе из 83 инженерных признаков, извлечённых из сырых PCAP-файлов или предварительно обработанных CSV-наборов данных. Признаки охватывают пять категорий, обобщённых в Таблице~\ref{tab:features}.

\begin{table}[htbp]
\centering
\caption{Конвейер инженерии 83 признаков, организованных по категориям.}
\label{tab:features}
\small
\begin{tabular}{p{3.5cm} p{1.5cm} p{8.5cm}}
\toprule
\textbf{Категория} & \textbf{Кол-во} & \textbf{Репрезентативные признаки} \\
\midrule
Статистика уровня потоков & 22 & Длительность, счётчики пакетов прямых/обратных, общее количество байтов (прям./обр.), длина пакетов (среднее, макс., мин., стд.), байт/с потока, пакетов/с потока, IAT потока (среднее, стд., макс., мин.) \\
\addlinespace
Распределения межпакетных интервалов & 18 & Прямой IAT (среднее, стд., мин., макс.), обратный IAT (среднее, стд., мин., макс.), статистики активного/неактивного времени, счётчики подпотоков прямых/обратных \\
\addlinespace
Флаговые индикаторы & 15 & Счётчики SYN, FIN, RST, PSH, ACK, URG, CWE, ECE, флаги PSH прямые/обратные, флаги URG прямые/обратные, счётчик флага FIN, счётчик флага SYN, счётчик флага RST \\
\addlinespace
Статистика полезной нагрузки & 16 & Размеры сегментов (мин., макс., среднее), длины заголовков (среднее, общее), начальные байты окна (прямые/обратные), энтропия байтов полезной нагрузки, дисперсия длины пакетов, средний размер пакета, средний размер сегмента \\
\addlinespace
Производные отношения & 12 & Байтов на пакет (прям./обр.), пакетов в секунду, соотношение нисх./восх., отношение среднего размера пакета, отношение IAT потока, массовая скорость (прямая/обратная), среднее/стд. активного периода, среднее/стд. неактивного периода \\
\midrule
\textbf{Всего} & \textbf{83} & \\
\bottomrule
\end{tabular}
\end{table}

Извлечение признаков выполняется общим модулем \texttt{FeatureExtractor} с использованием пофакторной z-стандартизации, подогнанной на обучающем распределении:
\begin{equation}
\tilde{f}_j = \frac{f_j - \mu_j^{\mathrm{train}}}{\sigma_j^{\mathrm{train}} + \varepsilon},
\label{eq:zscore}
\end{equation}
где $\mu_j^{\mathrm{train}}$ и $\sigma_j^{\mathrm{train}}$ --- среднее и стандартное отклонение признака $j$, вычисленные на обучающем множестве, а $\varepsilon=10^{-8}$ предотвращает деление на ноль. Наборы данных с менее чем 83 нативными признаками дополняются нулями до стандартной размерности; те, что имеют больше, проецируются через PCA с сохранением 95\% объяснённой дисперсии.


%%=====================================================================
\section{Предметно-ориентированная реализация методов}
\label{sec:method_instantiation}
%%=====================================================================

Данный раздел описывает, как каждый из семи методов Главы~\ref{ch:methods} реализован с параметрами, специфичными для кибербезопасности, демонстрируя отображение абстрактных математических объектов на конкретные сущности сетевой безопасности.


\subsection{М1: CT-TGNN для динамики графов сетевого трафика}

CT-TGNN~\cite{Anaedevha2025cais} реализован следующим образом. В реализации для сетевой безопасности множество узлов $\calV$ включает хосты и сервисы, а множество рёбер $\calE(t)$ регистрирует потоки. Признаки узлов $\bx_v(t)$ кодируют статистику уровня хоста (12 измерений). Признаки рёбер $\mathbf{e}_{uv}(t)$ захватывают 83-мерный вектор признаков потока из Таблицы~\ref{tab:features}.

Три параллельные ветви ОДУ из уравнения~\eqref{eq:ct_tgnn_multi} получают постоянные времени $\tau_1=1$\,с (переходные процессы уровня пакетов), $\tau_2=60$\,с (динамика сессий уровня потоков) и $\tau_3=3600$\,с (тенденции уровня кампаний). Это обеспечивает одновременное разрешение субсекундных всплесков сканирования портов и многочасовых кампаний развитых постоянных угроз. Архитектура использует четыре слоя передачи сообщений, скрытую размерность 256, размерность состояния 64 и адаптивное интегрирование DOPRI5 с допусками $10^{-5}$ (относительный) и $10^{-7}$ (абсолютный).

\begin{figure}[H]
\centering
\includegraphics[width=0.75\textwidth]{figures/fig_CT-TGNN_architecture.pdf}
\caption{Архитектура CT-TGNN, реализованная для динамики сетевого трафика.}
\label{fig:ct_tgnn_arch}
\end{figure}


\subsection{М2: FedLLM-API для совместной разведки угроз}

FedLLM-API~\cite{Anaedevha2026neurips} реализован для совместной разведки угроз. Каждый федеративный клиент соответствует центру операций безопасности, владеющему локальными наборами данных трафика. Византийско-устойчивая агрегация (Алгоритм~\ref{alg:fedllm}) защищает от скомпрометированных участников, отправляющих отравленные обновления градиентов. Дифференциальная приватность через локальную обрезку градиентов при пороге $C=1.0$ и гауссовский шум с масштабом $\sigma=0.8$ обеспечивает гарантию $(\epsilon{=}0.85,\delta{=}10^{-5})$ за 100 раундов коммуникации. Выравнивание ОТ Синкхорна с регуляризацией $\lambda=0.1$ корректирует распределительную гетерогенность между центрами. Вес ансамбля LLM $\lambda_{\mathrm{LLM}}=0.3$.

\begin{figure}[H]
\centering
\includegraphics[width=0.96\textwidth]{figures/fig_fedllm_api_architecture.pdf}
\caption{Архитектура FedLLM-API для совместной разведки угроз через организационные границы.}
\label{fig:fedllm_arch}
\end{figure}


\subsection{М3: TripleE-TGNN для многоуровневого мониторинга безопасности}

Граф сервисов $\calG_s(t)$ имеет узлы, представляющие микросервисы, и рёбра, представляющие API-вызовы. Граф трассировок $\calG_r(t)$ имеет узлы, представляющие распределённые трассировки. Граф хостов $\calG_n(t)$ имеет узлы, представляющие узлы инфраструктуры. Сила консолидации EWC $\lambda_{\mathrm{ewc}}=0.5$ балансирует пластичность и стабильность при эволюции топологии. Веса многозадачных потерь $\lambda_s=0.4$, $\lambda_r=0.3$, $\lambda_n=0.3$ с весом согласованности $\lambda_{\mathrm{con}}=0.1$.

\begin{figure}[H]
\centering
\includegraphics[width=0.96\textwidth]{figures/fig_triplee_tgnn_architecture.pdf}
\caption{Архитектура TripleE-TGNN для многоуровневого мониторинга безопасности в облачных средах.}
\label{fig:triplee_arch}
\end{figure}


\subsection{М4: MambaShield для обнаружения с адаптацией к дрейфу}

MambaShield~\cite{Anaedevha2026mamba} реализован для обнаружения с адаптацией к дрейфу. Последовательные записи потоков обрабатываются селективной моделью пространства состояний, входозависимые переходы которой адаптируются к текущему режиму трафика. Размерность состояния $d_s=64$, шаг дискретизации $\Delta=0.01$ и длина свёрточного ядра БПФ $L=1024$. Ансамбль учителей включает $M=5$ моделей, обученных на непересекающихся подмножествах данных. Учебный план начинается с доли отравления $\rho_0=0.05$ и увеличивается до $\rho_{\max}=0.3$. Дисперсия априорного PAC-Байеса $\sigma_P^2=1.0$.

\begin{figure}[H]
\centering
\includegraphics[width=0.75\textwidth]{figures/fig_mambashield_architecture.pdf}
\caption{Архитектура MambaShield для обнаружения вторжений с адаптацией к дрейфу и прогрессивной дистилляцией робастности.}
\label{fig:mamba_arch}
\end{figure}


\subsection{М5/М6: Оценка угроз с калиброванной неопределённостью}

Стохастический трансформер использует $M=50$ прямых проходов Монте-Карло при выводе для получения оценок эпистемической и алеаторной неопределённости. Предсказания с эпистемической неопределённостью, превышающей порог $\tau=0.3$, направляются аналитикам, снижая утомление от оповещений на 43\%. Модуль UC-HGP обеспечивает неопределённость на иерархических гауссовских процессах на временных масштабах, соответствующих ветвям CT-TGNN ($\tau_1$, $\tau_2$, $\tau_3$). Скорость вариационного обучения $\alpha_\phi=5\times10^{-4}$; скорость обучения возмущений $\alpha_\psi=10^{-3}$.

\begin{figure}[H]
\centering
\includegraphics[width=0.96\textwidth]{figures/fig_stochastic_transformer_architecture.pdf}
\caption{Стохастический трансформер с сортировкой на основе неопределённости для операций безопасности.}
\label{fig:stochastic_arch}
\end{figure}


\subsection{М7: Теоретико-игровая система для адаптивных противников}

Пространство действий защитника $\calC$ включает конфигурации обнаружения: покрытие мониторинга, пороги оповещений и политики реагирования. Пространство действий противника $\calA$ охватывает выбор цели, тип возмущения и стратегию уклонения. Скорость обучения мультипликативных весов $\eta=\sqrt{\log|\calC|/T}$. Штраф за неопределённость $\kappa=0.5$; порог направления эксперту $\tau=0.3$.

\begin{figure}[H]
\centering
\includegraphics[width=0.86\textwidth]{figures/fig_game_theoretical_architecture.pdf}
\caption{Теоретико-игровая система для адаптивных графовых противников в сетевой безопасности.}
\label{fig:game_arch}
\end{figure}


\subsection{Прямо совместимое обнаружение для постквантовых протоколов}

Протокол-специфичные кодировщики~\cite{Anaedevha2026edm,NIST2024,Regev2009} обрабатывают признаки тайминга для протоколов на основе решёток (Kyber/Dilithium~\cite{Avanzi2022}), структурные признаки для протоколов на основе кодов (McEliece/Niederreiter) и хэш-признаки (SPHINCS+). Слияние кросс-вниманием объединяет три представления. Бюджет состязательного обучения $\epsilon=0.03$ с возмущениями PGD-40.

\begin{figure}[H]
\centering
\includegraphics[width=0.96\textwidth]{figures/fig_pq_idps_architecture.pdf}
\caption{Прямо совместимая мультипротокольная архитектура обнаружения для постквантового трафика.}
\label{fig:pqidps_arch}
\end{figure}


%%=====================================================================
\section{Конфигурация гиперпараметров и оптимизация}
\label{sec:hyperparams}
%%=====================================================================

Таблица~\ref{tab:hyperparams} объединяет основные гиперпараметры всех методов в конфигурации для области сетевой безопасности. Все значения подобраны поиском по сетке на валидационном наборе CIC-IoT-2023 и подтверждены стабильными на оставшихся пяти наборах данных.

\begin{table}[htbp]
\centering
\caption{Сводная конфигурация гиперпараметров всех методов для области кибербезопасности.}
\label{tab:hyperparams}
\small
\begin{tabular}{p{3.5cm} p{2.5cm} p{2.5cm} p{5.0cm}}
\toprule
\textbf{Метод} & \textbf{Параметр} & \textbf{Значение} & \textbf{Обоснование} \\
\midrule
\multirow{4}{3.5cm}{М1: CT-TGNN}
& Скрыт. разм. & 256 & Ёмкость для 83-призн. входа \\
& Допуск ОДУ & $10^{-5}$ & Баланс точности/скорости \\
& Пост. времени & 1/60/3600\,с & Масшт. пакет/сессия/кампания \\
& Граница Липш. & Спектр. норма & Сертифиц. робастность \\
\addlinespace
\multirow{4}{3.5cm}{М2: FedLLM-API}
& Обрезка $C$ & 1.0 & Граница чувствит. ДП \\
& ДП $(\epsilon,\delta)$ & $(0.85, 10^{-5})$ & Строгая гарантия приватн. \\
& Рег. ОТ $\lambda$ & 0.1 & Сходимость Синкхорна \\
& Вес LLM & 0.3 & Баланс ансамбля \\
\addlinespace
\multirow{3}{3.5cm}{М3: TripleE-TGNN}
& EWC $\lambda$ & 0.5 & Предотвращение забывания \\
& Веса уровней & 0.4/0.3/0.3 & Акцент на сервисах \\
& Согласов. $\lambda_{\mathrm{con}}$ & 0.1 & Межуровн. согласованн. \\
\addlinespace
\multirow{3}{3.5cm}{М4: MambaShield}
& Разм. состояния & 64 & Ёмкость SSM \\
& Длина ядра & 1024 & Охват последовательн. \\
& Расп. отравления & 0.05--0.3 & Прогрессивн. закалка \\
\addlinespace
\multirow{3}{3.5cm}{М5/М6: Стох.\ трансф.}
& Выборки МК & 50 & Качество неопределённости \\
& Порог сортировки & 0.3 & Баланс ЛП/ЛО \\
& Калибровка & Температура & Минимизация ECE \\
\addlinespace
\multirow{2}{3.5cm}{М7: Сертификация}
& MWU $\eta$ & $\sqrt{\log|\calC|/T}$ & Оптимальный регрет \\
& Штраф неопр. $\kappa$ & 0.5 & Чувствит. к риску \\
\bottomrule
\end{tabular}
\end{table}

Все методы реализованы на PyTorch~2.1~\cite{Paszke2019} с CUDA~12.1. Обучение использует AdamW с начальной скоростью обучения $2\times10^{-4}$, затуханием весов $10^{-5}$ и косинусным отжигом. Обучение со смешанной точностью снижает потребление памяти на 40\%. Обрезка градиентов при норме 1.0 предотвращает нестабильность. Ранняя остановка с терпением 10 эпох мониторит потери валидации.


%%=====================================================================
\section{Интегрированная архитектура системы}
\label{sec:integrated_arch}
%%=====================================================================

Семь методов (М1--М7) интегрированы в единый конвейер обнаружения, поток данных которого проходит следующим образом. Сырой сетевой трафик поступает через модуль приёма данных, выполняющий разбор протоколов, извлечение потоков и вычисление 83 признаков. Модуль CT-TGNN строит темпоральный граф и вычисляет вложения узлов непрерывного времени. Модуль TripleE-TGNN поддерживает многоуровневые вложения для иерархических сред. Модуль MambaShield обрабатывает последовательные записи событий для оценки с адаптацией к дрейфу. Модули стохастического трансформера и UC-HGP обеспечивают калиброванные оценки неопределённости. Все выходы поступают в механизм принятия решений, применяющий теоретико-игровую стратегию из системы сертификации М7, формируя оповещения, классифицированные по степени серьёзности и уверенности. Модуль FedLLM-API обеспечивает периодическую синхронизацию моделей через организационные границы.

\begin{figure}[H]
\centering
\includegraphics[width=0.96\textwidth]{figures/Main_arch_v1b.pdf}
\caption{Интегрированный конвейер обнаружения, показывающий поток данных от сырого сетевого трафика через все семь методов к оповещениям, ранжированным по серьёзности и оценённым по уверенности. Модуль FedLLM-API синхронизирует модели через границы при дифференциальной приватности.}
\label{fig:u_idps_arch}
\end{figure}

Система предназначена для развёртывания как плагин для существующих систем СОПВ, таких как Snort~3~\cite{Roesch1999} и Suricata, принимая их потоки оповещений в качестве дополнительного входа и формируя обогащённые выходы с оценкой уверенности. Полная программная реализация описана в Главе~\ref{ch:platform}, публично архивирована под DOI 10.5281/zenodo.19129512~\cite{Anaedevha2026robustidps}.


%%=====================================================================
\section{Итоги главы}
\label{sec:ch3_summary}
%%=====================================================================

В настоящей главе продемонстрировано, как предметно-независимые математические методы Главы~\ref{ch:methods} отображаются на сущности и ограничения обнаружения сетевых вторжений. Ключевые вклады данной главы: (1)~формальный конвейер преобразования данных в граф, конвертирующий сырой сетевой трафик PCAP/CSV во временные графовые структуры с признаками узлов (12 измерений) и рёбер (83 измерения); (2)~унифицированная таксономия классификации из 34 классов по 7 категориям атак, расширяемая через классификацию zero-shot LLM; (3)~комплексный конвейер инженерии 83 признаков, охватывающий статистику потоков, распределения межпакетных интервалов, флаговые индикаторы, статистику полезной нагрузки и производные отношения; (4)~предметно-ориентированная реализация всех семи методов с гиперпараметрами, постоянными времени и архитектурными конфигурациями, адаптированными для кибербезопасности; и (5)~интегрированный конвейер обнаружения, демонстрирующий сквозной поток от сырого трафика до оповещений с оценкой уверенности. В следующей главе представлена полная экспериментальная программа, валидирующая эти предметно-ориентированные реализации на шести бенчмарк-наборах данных общим объёмом 84,2 миллиона размеченных образцов.


